\documentclass[conference]{IEEEtran}
\usepackage{amsmath,amssymb,amsfonts}
\usepackage{booktabs}
\usepackage{graphicx}
\usepackage{textcomp}
\usepackage{xcolor}
\usepackage{hyperref}

\begin{document}

\title{Activations, Not Weights: Orthogonal Preconditioning for\\
Low-Bit Quantization of Vision--Language--Action Policies}

\author{
\IEEEauthorblockN{Justin Williams, Kishor Datta Gupta, Roy George, Mrinmoy Sarkar}
\IEEEauthorblockA{Clark Atlanta University \quad \texttt{justinwilliamstech693@gmail.com}}
}

\maketitle

\begin{abstract}
Vision--Language--Action (VLA) policies built on 7B-parameter backbones are accurate but too large
for on-board robot inference. We present a comprehensive, closed-loop study of post-training
quantization for an OpenVLA-OFT policy on the LIBERO benchmark, distilled into one actionable
principle: \emph{quantization wants activation energy spread uniformly across channels.}
(i)~\textbf{Weights are easy}---naive per-channel quantization is near-lossless to 4 bits
(90.75\% vs.\ a 90.25\% bf16 baseline). (ii)~\textbf{Activations are the wall}---naive 4-bit
activation quantization collapses the policy to 0\% because a few channels carry magnitudes up to
$10^5\times$ the median. (iii)~A single orthogonal \emph{incoherence rotation} restores W4A4 to
88.8--90.5\% (98--100\% retention) by cutting the median activation outlier ratio from 18.0 to 1.6.
We benchmark \textbf{16 transforms} and distill a single taxonomy---global dense mixers work,
localized or energy-concentrating transforms do not (PCA is actively harmful). \textbf{DCT} is the
deployment pick (data-free, $O(n\log n)$, no power-of-two fallback). Component-wise allocation
(vision INT8 $+$ DCT-W3A8 backbone $+$ head INT8) yields \textbf{74\% compression at no measurable
accuracy loss}. We confirm the result on \textbf{real Jetson AGX Orin} hardware: W4A4+DCT achieves
86.0\% on LIBERO-Spatial ($-2.2$\,pt vs.\ A100 bf16; statistically indistinguishable at 97.5\%
retention, 73\% smaller). Code, harness, and leaderboard are released.
\end{abstract}

\begin{IEEEkeywords}
VLA quantization, activation outliers, incoherence rotation, DCT, edge deployment, Jetson Orin, LIBERO
\end{IEEEkeywords}

\section{Introduction}

Vision--Language--Action models turn large vision--language backbones into robot policies that map
an image and a language instruction to low-level actions~\cite{openvla,openvlaoft}. Their accuracy
rests on 7B-parameter backbones, but deploying such models on embedded robot compute demands
aggressive compression. Post-training quantization (PTQ) is the natural tool---yet most PTQ
research optimizes language-model perplexity or open-loop token accuracy~\cite{gptq,awq,smoothquant,quarot,spinquant},
not the \emph{closed-loop} task success where per-step error compounds over long horizons.

This paper answers three coupled questions: \emph{where} does quantization hurt a VLA policy,
\emph{what} is the cheapest fix, and \emph{how far} can one push before accuracy breaks. Our
contributions:
\begin{enumerate}
  \item \textbf{Diagnosis.} The bottleneck is activation outliers, not weights: ``massive
        activation'' channels up to $10^5\times$ the median collapse naive W4A4 to 0\% success
        (Sec.~\ref{sec:mech}).
  \item \textbf{Fix.} A single orthogonal incoherence rotation enables W4A4 at 98--100\%
        retention (Sec.~\ref{sec:wall}).
  \item \textbf{Benchmark.} 16 transforms ranked by a seconds-long proxy; one taxonomy predicts
        every result; DCT is the deployment pick (Sec.~\ref{sec:zoo}).
  \item \textbf{Deployment.} Component-wise allocation gives 74\% compression at no measurable
        accuracy loss; confirmed on real Jetson AGX Orin at 97.5\% retention
        (Sec.~\ref{sec:component}--\ref{sec:device}).
\end{enumerate}

\section{Background and Related Work}

\textbf{VLA policies.}
OpenVLA~\cite{openvla} and OpenVLA-OFT~\cite{openvlaoft} couple a vision encoder with a
Llama-2-7B backbone and an action head; OFT regresses a continuous 8-action chunk from the
backbone's last-layer hidden states.

\textbf{PTQ and outliers.}
Weight-only PTQ (GPTQ~\cite{gptq}, AWQ~\cite{awq}) is mature; activation quantization is
bottlenecked by outlier channels. SmoothQuant~\cite{smoothquant} migrates scale from activations
to weights; incoherence processing (QuaRot~\cite{quarot}, SpinQuant~\cite{spinquant},
QuIP~\cite{quip}) inserts an orthogonal rotation $Q$ so $Wx = (WQ^\top)(Qx)$ holds exactly while
the rotated operands become near-Gaussian.

\textbf{Concurrent VLA quantization.}
$\Omega$-QVLA~\cite{omegaqvla} pushes diffusion-head VLAs ($\pi_{0.5}$, GR00T) to W4A4 (98.0\%
retention) with a composite SVD--Hadamard rotation. DyQ-VLA~\cite{dyqvla} reaches 99.5\% retention
via temporal-dynamic bit-width switching (1.43$\times$ measured speedup). Both target diffusion
heads; we target a \emph{non-diffusion} L1-regression OFT head where per-step DiT scaling does not
apply, and we benchmark 16 transforms rather than picking one.

\section{Method}
\label{sec:method}

We quantize only the Llama backbone linears (84\% of parameters); the vision tower, projector,
lm\_head, embeddings, and regression head stay bf16. All experiments use \emph{simulated}
(fake) quantization---quantize then dequantize in floating point---so that closed-loop success
isolates the algorithm from kernel artifacts.

\textbf{Quantizers.}
Weights: symmetric per-output-channel, $\hat{w} = \mathrm{round}(w/s)s$, $s = \max|w|/(2^{b-1}-1)$.
Activations: symmetric per-token at $a$ bits. A configuration is denoted W$b$A$a$.

\textbf{Preconditioning.}
For each backbone linear we insert an invertible preconditioner $M$ on the input axis:
\begin{equation}
  y = Wx = (WM^{-1})(Mx)
\end{equation}
For orthogonal $M$, $M^{-1}=M^\top$ and error norms are basis-invariant.

\textbf{Proxy.}
We score transforms offline by the median INT4 per-tensor activation error after applying $M$:
\begin{equation}
  \mathrm{A4err}(x) = \frac{\mathbb{E}|\hat{x} - x|}{\mathbb{E}|x|}, \quad
  \hat{x} = \mathrm{dequant}_4^{\mathrm{per\text{-}tensor}}(Mx)
\end{equation}
over a single 32-frame calibration pass---ranking all 16 transforms in seconds.

\section{Experiments}

\textbf{Setup.} OpenVLA-OFT (bf16 success 90.25\%); closed-loop on all four LIBERO
suites~\cite{libero}, 10 tasks $\times$ 10 rollouts $= 400$ rollouts per configuration.
Binomial CI at 400 rollouts $\approx \pm 2.5\%$.

\subsection{Weights Are Easy; Activations Are the Wall}
\label{sec:wall}

Naive weight quantization is near-lossless to 4 bits, but naive W4A4 collapses to 0\%.
A single rotation restores W4A4 to near-baseline; multiple global rotations cluster at
88.8--90.5\% (Table~\ref{tab:main}).

\begin{table}[h]\centering
\caption{Full closed-loop evaluation (400 rollouts) vs.\ bf16 baseline (90.25\%).}
\label{tab:main}
\begin{tabular}{lccc}
\toprule
Configuration & Acts & Success & Retention \\
\midrule
bf16 baseline & 16-bit & 90.25\% & --- \\
Naive W4 (per-channel) & 16-bit & 90.75\% & 100.6\% \\
Rotated W4 & 16-bit & 89.0\% & 98.6\% \\
\textbf{Naive W4A4} & \textbf{4-bit} & \textbf{0.0\%} & \textbf{0\%} \\
\textbf{Rot.\ W4A4 (Hadamard)} & \textbf{4-bit} & \textbf{88.8\%} & \textbf{98.4\%} \\
Rot.\ W4A4 (random-orth.) & 4-bit & 90.5\% & 100.3\% \\
Rot.\ W3A8 & 8-bit & 88.8\% & 98.4\% \\
\bottomrule
\end{tabular}
\end{table}

The cost concentrates in long-horizon tasks (Table~\ref{tab:persuite}), consistent with
per-step error accumulating over long rollouts.

\begin{table}[h]\centering
\caption{Per-suite success (\%) at W4A4 vs.\ bf16.}
\label{tab:persuite}
\begin{tabular}{lcccc}
\toprule
Suite & bf16 & Rot.\ W4A4 & Rot.\ W4 & Rot.\ W3A8 \\
\midrule
spatial & 87 & 87 & 88 & 87 \\
object  & 99 & 97 & 99 & 98 \\
goal    & 86 & 88 & 86 & 84 \\
long    & 89 & 83 & 83 & 86 \\
\midrule
avg & 90.25 & 88.8 & 89.0 & 88.8 \\
\bottomrule
\end{tabular}
\end{table}

\subsection{Mechanism: Rotation Crushes Activation Outliers}
\label{sec:mech}

Rotation cuts the median outlier ratio $11\times$ and the worst by four orders of magnitude
(Table~\ref{tab:mech}). The worst single layer (\texttt{layers.1.mlp.down\_proj}) carries a
channel $\sim 10^5\times$ the median---a ``massive activation''---which alone sets the per-tensor
scale and zeroes the rest, explaining the 0\% naive-W4A4 result.

\begin{table}[h]\centering
\caption{Activation conditioning before vs.\ after rotation (median over 224 backbone linears).}
\label{tab:mech}
\begin{tabular}{lccc}
\toprule
Metric & Pre & Post & Factor \\
\midrule
median outlier ratio & 18.0 & 1.6 & $11\times$ \\
max outlier ratio & 105,928 & 6.2 & $17{,}000\times$ \\
median A4 per-tensor error & 96.5\% & 26.1\% & $3.7\times$ \\
\bottomrule
\end{tabular}
\end{table}

\subsection{Which Transform? A Benchmark of 16 with a Taxonomy}
\label{sec:zoo}

One principle organizes all 16 results: \emph{spread energy, do not concentrate it.}
Table~\ref{tab:rank} reports the proxy ranking and closed-loop validation.

\begin{table}[h]\centering
\caption{Transform ranking by proxy (lower = better; identity $= 96.3\%$) with full-eval W4A4.}
\label{tab:rank}
\begin{tabular}{rlccc}
\toprule
\# & Transform & A4 err & Full eval & Category \\
\midrule
1 & SRHT & 26.3\% & 89.8\% & global mixer \\
2 & DST & 27.0\% & --- & global mixer \\
3 & \textbf{DCT} & \textbf{27.3\%} & \textbf{89.8\%} & global mixer \\
4 & Hadamard & 27.4\% & 88.8\% & global mixer \\
5 & Walsh & 27.4\% & --- & global mixer \\
6 & Hartley & 27.5\% & --- & global mixer \\
7 & rand-orth & 28.2\% & 90.5\% & global mixer \\
8 & polar & 70.7\% & --- & partial \\
9 & butterfly & 71.0\% & --- & partial \\
10 & \textbf{PCA} & \textbf{81.7\%} & --- & \textbf{harmful (concentrates)} \\
11--16 & others & 90--97\% & --- & localized / partial \\
\bottomrule
\end{tabular}
\end{table}

Global dense mixers blend every channel into every output $\Rightarrow$ work ($\sim$27\%).
PCA concentrates energy into top components---manufacturing a giant outlier. The proxy correctly
predicts ordering; DCT is the deployment pick: no power-of-two fallback, full-eval-confirmed at
89.8\% (99.5\% retention), data-free and $O(n\log n)$.

\subsection{Component-Wise Allocation: Vision Is Free, Head Is the Floor}
\label{sec:component}

Treating the policy as three quantizable components---fused vision tower, Llama backbone,
L1-regression head---reveals a clean sensitivity ordering (full eval, LIBERO-Spatial $10\times10$;
retention vs.\ bf16 86.0\%). The vision tower is \emph{free}: INT8 and FP8-E4M3 are both
lossless. The action head is the \emph{floor}: INT8 is safe, INT4 collapses to 70\%, INT3 to 50\%.
The deployment operating point is INT8~$+$~DCT-W4A4~$+$~INT8 (Table~\ref{tab:component}).

\begin{table}[h]\centering
\caption{Component-wise Pareto frontier (LIBERO-Spatial $10\times10$; retention vs.\ 86.0\%).}
\label{tab:component}
\begin{tabular}{lccc}
\toprule
Vision / Backbone / Head & Success & Size & Retention \\
\midrule
bf16 / bf16 / bf16 & 86.0\% & 15.4\,GB & 100\% \\
INT8 / INT8 / INT8 & 85.0\% & 8.0\,GB & 98.8\% \\
bf16 / DCT-W4A4 / bf16 & 84.0\% & 5.7\,GB & 97.7\% \\
INT8 / DCT-W4A4 / INT8 & 82.0\% & 4.8\,GB & 95.3\% \\
\textbf{INT8 / DCT-W3A8 / INT8} & \textbf{84.0\%} & \textbf{4.0\,GB} & \textbf{97.7\%} \\
\bottomrule
\end{tabular}
\end{table}

\textbf{All-suite headline.}
Across all four LIBERO suites (400 rollouts, bf16 baseline 88.0\%), \textbf{INT8 $+$ DCT-W3A8
$+$ INT8 achieves 89.5\%}---a \textbf{74\% smaller model at no measurable accuracy loss}
(within the $\pm2.5\%$ CI at 400 rollouts). Notably W3A8 matches/exceeds W4A4 at smaller size:
once the DCT rotation conditions the activations, weight precision is cheap and activation
precision is what remains.

\subsection{On-Device Realization}
\label{sec:device}

We deploy the quantized policy on a Jetson AGX Orin ($\sim$\$2k, $\sim$60\,W TDP). A single
backbone query takes 330\,ms; because OFT emits 8-action chunks, the effective control rate
is $\approx$24\,Hz---within the regime for LIBERO-style manipulation.

\begin{table}[t]\centering
\caption{W4A4+DCT on Jetson AGX Orin vs.\ A100 bf16 reference (220 steps, $10\times10$ rollouts).}
\label{tab:device}
\begin{tabular}{lcc}
\toprule
Metric & A100 bf16 & Orin W4A4+DCT \\
\midrule
Closed-loop success & 88.2\% & \textbf{86.0\%} \\
Retention vs.\ A100 & 100\% & \textbf{97.5\%} \\
Accuracy gap & --- & $-2.2$\,pt \\
Query latency & 53\,ms & 330\,ms \\
Effective control rate & $>$100\,Hz & $\approx$24\,Hz \\
Peak memory & 15.4\,GB & $\sim$4.1\,GB \\
Energy per query & --- & [tegrastats pending] \\
\bottomrule
\end{tabular}
\end{table}

The $-2.2$\,pt gap is within the binomial $\pm4\%$ CI at 100 rollouts---statistically
indistinguishable from the A100 result. This is the first \emph{real-hardware closed-loop
confirmation} of quantized L1-regression VLA task success on edge compute; prior VLA-quant
work (Omega-QVLA, DyQ-VLA) reports simulation-only accuracy results.

\section{Discussion and Limitations}

The remaining gap to a full-conference submission is energy/power measurement via
\texttt{tegrastats} and at least one physical robot rollout. We study one VLA and one benchmark;
generalization to other architectures and tasks is open. Numbers are single-seed (A100 sweeps:
$\pm2.5\%$ CI at 400 rollouts; Orin comparison: $\pm4\%$ CI at 100 rollouts). The proxy predicts
ordering, not absolute magnitude.

\section{Conclusion}

For closed-loop VLA control, activation outliers---not weights---are the quantization bottleneck,
and a single orthogonal rotation removes them, enabling W4A4 at 98--100\% retention and 2.7$\times$
compression. Across 16 transforms, one principle predicts everything---spread energy, do not
concentrate it; global mixers tie at $\sim$90\%, DCT is the deployment pick, PCA is a cautionary
tale. Component-wise allocation yields 74\% compression at no measurable accuracy loss, confirmed
on real Jetson AGX Orin at 97.5\% retention. Code, harness, and leaderboard are released.

\begin{thebibliography}{99}

\bibitem{openvla}
  A.~Kim et al., ``OpenVLA: An Open-Source Vision-Language-Action Model,''
  \textit{arXiv:2406.09246}, 2024.

\bibitem{openvlaoft}
  M.~Kim et al., ``Fine-Tuning Vision-Language-Action Models (OpenVLA-OFT),''
  \textit{arXiv:2409.12191}, 2025.

\bibitem{libero}
  B.~Liu et al., ``LIBERO: Benchmarking Knowledge Transfer for Lifelong Robot Learning,'' 2023.

\bibitem{gptq}
  E.~Frantar et al., ``GPTQ: Accurate Post-Training Quantization for GPT,'' 2022.

\bibitem{awq}
  J.~Lin et al., ``AWQ: Activation-aware Weight Quantization for LLM Compression,'' 2023.

\bibitem{smoothquant}
  G.~Xiao et al., ``SmoothQuant: Accurate and Efficient PTQ for LLMs,'' 2022.

\bibitem{quarot}
  S.~Ashkboos et al., ``QuaRot: Outlier-Free 4-Bit Inference in Rotated LLMs,'' 2024.

\bibitem{spinquant}
  Z.~Liu et al., ``SpinQuant: LLM Quantization with Learned Rotations,'' 2024.

\bibitem{quip}
  J.~Chee et al., ``QuIP/QuIP\#: 2-Bit Quantization with Incoherence Processing,'' 2023--2024.

\bibitem{omegaqvla}
  ``$\Omega$-QVLA: Robust Quantization for VLAs via Composite Rotation and Per-step Scaling,''
  2026.

\bibitem{dyqvla}
  ``DyQ-VLA: Temporal-Dynamic-Aware Quantization for Embodied VLAs,''
  \textit{arXiv:2603.07904}, 2026.

\end{thebibliography}

\end{document}
